% !TEX root = ../Main.tex
\chapter{平面向量基础复习}

高中阶段的空间几何计算以\textbf{向量法}为主。为了顺利过渡到空间，必须先把\textbf{平面向量}的几条核心结论回忆清楚——\textbf{共线、基本定理、点乘、模长、夹角、平行、垂直}。这些结论在空间中只需\textbf{升一维}即可直接使用。

\section{向量的坐标与共线}

\state{平面向量坐标}{
    平面向量 $\vec a = (x, y)$ 表示 $\vec a = x \vec i + y \vec j$，其中 $\vec i, \vec j$ 是\textbf{标准基}。

    \textbf{两向量共线判据}（平行 $\vec a \parallel \vec b$）：
    \[
    \vec a \parallel \vec b \iff \exists \lambda \in \mathbb R,\ \vec b = \lambda \vec a \iff x_1 y_2 - x_2 y_1 = 0.
    \]
}

\state{三点共线与 $\lambda + \mu = 1$}{
    设 $A, B, C$ 为平面上三个点，$O$ 为任意参考点。若
    \[
    \vec{OC} = \lambda \vec{OA} + \mu \vec{OB},
    \]
    则
    \[
    \boxed{A, B, C \text{ 三点共线} \iff \lambda + \mu = 1.}
    \]

    \textbf{证}：三点共线 $\iff \vec{AC} = t \vec{AB}$，即 $\vec{OC} - \vec{OA} = t(\vec{OB} - \vec{OA})$，整理得 $\vec{OC} = (1 - t) \vec{OA} + t \vec{OB}$，系数和为 $1$。
}

\section{平面向量基本定理}

\state{平面向量基本定理}{
    若 $\vec e_1, \vec e_2$ 为\textbf{不共线}的两个向量（即平面上的一组基），则平面内\textbf{任意}向量 $\vec a$ 都能\textbf{唯一}表示成
    \[
    \vec a = \lambda_1 \vec e_1 + \lambda_2 \vec e_2.
    \]
}

\rmkb{
    注意：\textbf{非共线}是关键——共线的两个向量只能张成一条直线，而非整个平面。
}

\section{数量积（点乘）}

\state{点乘的几何定义}{
    两向量 $\vec a, \vec b$ 夹角为 $\theta$，则
    \[
    \boxed{\vec a \cdot \vec b = |\vec a|\, |\vec b| \cos \theta.}
    \]

    \textbf{几何意义}：
    \[
    \vec a \cdot \vec b = \underbrace{|\vec a| \cos \theta}_{\vec a \text{ 在 } \vec b \text{ 方向上的投影}} \times |\vec b|.
    \]
    即 \textbf{"投影 $\times$ 模长"}。当 $\vec b$ 为单位向量 $\hat{\vec b}$ 时，$\vec a \cdot \hat{\vec b}$ 就是 $\vec a$ 在 $\hat{\vec b}$ 方向上的\textbf{投影长度}。
}

\state{点乘的坐标形式}{
    若 $\vec a = (x_1, y_1),\ \vec b = (x_2, y_2)$，则
    \[
    \boxed{\vec a \cdot \vec b = x_1 x_2 + y_1 y_2.}
    \]
}

\section{模长、夹角、平行、垂直}

\state{模长}{
    $\vec a = (x, y) \Rightarrow |\vec a| = \sqrt{x^2 + y^2}$。特别地 $|\vec a|^2 = \vec a \cdot \vec a$。
}

\state{夹角公式}{
    $\vec a, \vec b$ 的夹角 $\theta$ 满足
    \[
    \cos \theta = \frac{\vec a \cdot \vec b}{|\vec a|\, |\vec b|} = \frac{x_1 x_2 + y_1 y_2}{\sqrt{x_1^2 + y_1^2}\, \sqrt{x_2^2 + y_2^2}}.
    \]
}

\state{平行与垂直判据}{
    \begin{itemize}
        \item \textbf{平行}（$\vec a \parallel \vec b$）$\iff x_1 y_2 - x_2 y_1 = 0$；
        \item \textbf{垂直}（$\vec a \perp \vec b$）$\iff \vec a \cdot \vec b = 0 \iff x_1 x_2 + y_1 y_2 = 0$。
    \end{itemize}
}

\section{升维：从平面到空间}

\rmkb{
    \textbf{以上所有结论在空间中只需加一维即可使用}：
    \begin{itemize}
        \item 向量 $\vec a = (x, y, z)$；
        \item 点乘 $\vec a \cdot \vec b = x_1 x_2 + y_1 y_2 + z_1 z_2$；
        \item 模长 $|\vec a| = \sqrt{x^2 + y^2 + z^2}$；
        \item 垂直判据 $\vec a \perp \vec b \iff \vec a \cdot \vec b = 0$；
        \item 夹角 $\cos \theta = \dfrac{\vec a \cdot \vec b}{|\vec a|\, |\vec b|}$。
    \end{itemize}
    \textbf{唯一的新概念}：空间向量的\textbf{平行}不再由"叉积"单条件判定，而需 $\exists \lambda$ 使 $\vec b = \lambda \vec a$（坐标对应成比例）。
}
